
@@M0@@
Les
foncteurs @@M1@@, @@M2@@ et @@M3@@ transforment modules libres (resp. localement libres) en modules libres (resp. localement libres).

De façon générale, considérons des sections @@M4@@ de @@M5@@ et un multifoncteur (covariant, pour fixer les idées)
@@M6@@
à @@M7@@ variables, où @@M8@@ désigne la catégorie fibrée sur @@M9@@ des @@M10@@-modules localement libres de rang @@M11@@. Soit
@@M12@@
et
@@M13@@
Il s'ensuit immédiatement que @@M14@@ définit une représentation
@@M15@@
Il est clair que si @@M16@@ est lui-même localement libre de rang @@M17@@, pour tout foncteur composé,
@@M18@@
on a
@@M19@@

Réciproquement, étant donné la représentation @@M20@@, on récupère @@M21@@ (pour toute suite
